\begin{abstract}
If we read 15 documents from a corpus, how much of the corpus have we seen? Researchers who study text read samples because they cannot analyze every document. A survey sample comes with a formula for its size and a margin of error, and a sample of text comes with neither. We measure what a sample has missed as the percentage of the corpus with no document of the same content in the sample, where we define two documents as the same content when their embeddings lie closer than a distance the researcher chooses. We estimate the missed percentage from the sample alone with a classical estimator built on what we call the isolation rate, the percentage of sampled documents with no other sampled document of the same content. To plan the size of a sample, we show that on most corpora we tested a larger sample and a looser definition of the same content trade against each other at a fixed rate. The rate is set by the intrinsic dimension of the corpus, the number of independent ways in which its documents differ in meaning. This rate gives a sample size formula. In 77 of 83 tests, each a combination of corpus, target and distance, the sample actually needed fell between two thirds and one and a half times the formula's answer. A pilot of 2{,}000 documents shows in advance whether the formula can be trusted for a given corpus. Two large language models judged pairs of documents to tell which distance means the same specific content and which means the same topic.
\end{abstract}
